% Solution: Gibbons (1992) Problem 2.1 "The Rotten Kid Theorem". Input by \solitem.
% Reused from lista-reais-lectures-01-04-solutions/item18.tex, deepened (exact worked example,
% role of each assumption, common mistakes). Check: check_B_L03.py.
\begin{exsol}{Gibbons Problem 2.1}{130}{Advanced}{5}
\asked Prove Becker's Rotten Kid theorem: when a child's action $A$ produces incomes
$I_C(A)$, $I_P(A)$ and an altruistic parent then chooses a bequest $B$, the backwards-induction
outcome has the child choosing the $A$ that maximizes $I_C(A) + I_P(A)$.

\method Backward induction in a two-stage game of perfect information with continuous actions
(Lecture~3, slides~29--30; Gibbons \S2.1.A, pp.~57--61; Tadelis \S8.3.2): solve the parent's
problem for every pair of incomes, then the child's problem given the parent's reply. The key step
is that the parent's reply depends on the incomes only through their sum, and raises the child's
consumption when the sum rises.

\solution
\textbf{Notation and assumptions.} $A \ge 0$ the child's action; $I_C, I_P$ strictly concave with
maximizers $A_C, A_P > 0$; $B \in \mathbb{R}$ the bequest; $U, V$ increasing and strictly concave;
$k > 0$. Write $Y = I_C + I_P$ (family income) and $x = I_C + B$ (the child's consumption), so the
parent consumes $I_P - B = Y - x$. As is standard for this problem, $U$ and $V$ are
differentiable and the parent's problem has a solution.

\textbf{Step 1: the parent's problem depends only on $Y$.} Given $(I_C, I_P)$ the parent solves
\[
  \max_{B \in \mathbb{R}}\ V(I_P - B) + k\,U(I_C + B).
\]
Because $B$ is unrestricted in sign, choosing $B$ is the same as choosing $x = I_C + B \in \mathbb{R}$,
and the problem becomes
\begin{equation}
  \max_{x \in \mathbb{R}}\ W(x; Y) = V(Y - x) + k\,U(x). \label{eq:gibbons-2-1-W}
\end{equation}
$W$ is strictly concave in $x$ (a sum of strictly concave functions), so the maximizer $x^*(Y)$ is
unique and satisfies the first-order condition
\begin{equation}
  V'(Y - x^*) = k\,U'(x^*). \label{eq:gibbons-2-1-foc}
\end{equation}
The bequest is $B^*(I_C, I_P) = x^*(I_C + I_P) - I_C$: it depends on both incomes, but the child's
consumption depends only on their sum.

\textbf{Step 2: $x^*$ is strictly increasing in $Y$.} Take $Y' > Y$ and suppose
$x' = x^*(Y') \le x = x^*(Y)$. Then $Y' - x' > Y - x$. Strict concavity makes $V'$ strictly
decreasing and $U'$ decreasing, so
\[
  V'(Y' - x') < V'(Y - x) \qquad\text{and}\qquad k\,U'(x') \ge k\,U'(x).
\]
Applying \eqref{eq:gibbons-2-1-foc} at $Y$ and at $Y'$:
\[
  V'(Y' - x') = k\,U'(x') \;\ge\; k\,U'(x) = V'(Y - x) \;>\; V'(Y' - x'),
\]
a contradiction. Hence $x^*(Y') > x^*(Y)$. (With $U, V$ twice differentiable,
$dx^*/dY = V''/(V'' + kU'') \in (0,1)$: the parent passes part, not all, of an income gain to the
child.)

\textbf{Step 3: the child's problem.} Anticipating $B^*$, the child chooses $A \ge 0$ to maximize
\[
  U\bigl(I_C(A) + B^*(I_C(A), I_P(A))\bigr) = U\bigl(x^*(I_C(A) + I_P(A))\bigr) = U\bigl(x^*(Y(A))\bigr).
\]
$U \circ x^*$ is strictly increasing (Steps~1--2 and $U$ increasing), so $A$ maximizes the child's
payoff if and only if it maximizes $Y(A) = I_C(A) + I_P(A)$.

\textbf{Step 4: the maximizer exists and is unique.} $Y$ is strictly concave. For
$A < \min\{A_C, A_P\}$ both incomes are increasing, and for $A > \max\{A_C, A_P\}$ both are
decreasing, so $Y$ attains its maximum on the compact interval $[\min\{A_C,A_P\}, \max\{A_C,A_P\}]
\subset (0, \infty)$, at a unique point $A^*$. The backwards-induction outcome is
\[
  \bigl(A^*,\ B^*(I_C(A^*), I_P(A^*))\bigr), \qquad A^* = \arg\max_{A \ge 0}\ [\,I_C(A) + I_P(A)\,].
  \qquad\square
\]

\textbf{Worked example (exact).} $U = V = \ln$, $k = 2$, $I_C = 4A - A^2$ ($A_C = 2$),
$I_P = 10A - A^2$ ($A_P = 5$). The parent's condition \eqref{eq:gibbons-2-1-foc} reads
$\frac{1}{Y - x} = \frac{2}{x}$, so $x^*(Y) = \frac{2Y}{3}$ (in general $x^* = \frac{kY}{1+k}$).
The child maximizes $\ln\bigl(\frac23(14A - 2A^2)\bigr)$: $14 - 4A = 0$, $A^* = \frac72$, between
$A_C = 2$ and $A_P = 5$. There $Y = \frac{49}{2}$, $x^* = \frac{49}{3}$, $I_C = \frac74$, so the
bequest is $B^* = \frac{49}{3} - \frac74 = \frac{175}{12}$. Had the child chosen its own income
maximizer $A_C = 2$, family income would be $Y(2) = 20 < \frac{49}{2}$ and its consumption only
$\frac{40}{3} < \frac{49}{3}$.

\interp The child is selfish but faces a parent who ``taxes'' any loss of family income: if the
child raised $I_C$ at a larger cost to $I_P$, the parent, who cares about the split, would cut the
bequest by more than the child gained. Altruism on one side, plus the parent's last move, aligns
the first mover's incentives with joint surplus. Each assumption has a role:
\begin{itemize}
  \item \emph{the parent moves last}: it is the parent's credible, sequentially rational reply that
  disciplines the child (Lecture~4's theme: what matters is what a player will actually do when the
  moment comes);
  \item \emph{$B$ unrestricted}: if $B \ge 0$ were imposed and binding, the child's consumption would
  be $I_C$ alone over a range and Step~1 fails;
  \item \emph{the child moves first and cannot commit the parent}: if the bequest were fixed before
  $A$, the child would simply choose $A_C$.
\end{itemize}

\verify \texttt{check\_B\_L03.py} checks with sympy that $dx^*/dY = V''/(V''+kU'') \in (0,1)$ for
several negative $(V'', U'')$ and positive $k$; derives $x^* = kY/(1+k)$ for log utility (second
derivative negative) and $x^* = (Y + \ln k)/2$ for CARA utility $-e^{-x}$; solves the child's problem
in the log example ($A^* = \frac72$), confirms on an exact grid of step $\frac1{20}$ that the
maximizer of $\frac23 Y(A)$ is $\frac72$ while $I_C$ alone peaks at $2$, and recomputes
$B^* = \frac{175}{12}$. Output line: \texttt{gibbons-2-1: x*(Y) increasing; child picks A* = 7/2 =
argmax(I\_C + I\_P), not A\_C = 2}. Gibbons prints no answer.

\mistakes
\begin{itemize}
  \item Assuming the child maximizes $I_C$: its payoff depends on the bequest, which reacts to $A$.
  \item Treating the bequest as a fixed number instead of a function $B^*(I_C, I_P)$ (a strategy of
  the parent in the second stage).
  \item Proving only that $x^*$ is weakly increasing, or skipping existence/uniqueness of $A^*$;
  the strict monotonicity in Step~2 is what makes ``maximize $U(x^*(Y))$'' equivalent to
  ``maximize $Y$''.
\end{itemize}

\answer{The parent's best bequest gives the child consumption $x^*(I_C + I_P)$ that solves
$V'(Y - x) = kU'(x)$; it depends only on family income $Y$ and is strictly increasing in it. So the
child's payoff $U(x^*(Y(A)))$ is a strictly increasing function of $Y(A)$, and the child chooses
$A^* = \arg\max_A [I_C(A) + I_P(A)]$, which exists and is unique by strict concavity.}
\end{exsol}
